\documentclass[12]{amsart}
\usepackage{amsmath, amssymb, amsthm}
\usepackage{geometry}
\usepackage{mathrsfs}
\usepackage{lmodern}
\usepackage{graphicx}
\usepackage{epigraph}
\usepackage{fancyhdr}
\geometry{margin=1in}
\pagestyle{fancy}
\lhead{Spectral Fixed Point Theory}
\rhead{\thepage}
\fancyfoot[C]{}        % Clear center footer (where default page number appears)
\fancyfoot[L]{}        % Optional: clear left footer
\fancyfoot[R]{}        % Optional: clear right footer
\renewcommand{\headrulewidth}{0.4pt}  % keep or customize header line
\renewcommand{\footrulewidth}{0pt}    
% Define theorem style (optional)
\theoremstyle{plain}
\newtheorem{theorem}{Theorem}[section]
\newtheorem{lemma}[theorem]{Lemma}
\newtheorem{definition}[theorem]{Definition}
\begin{document}
\begin{document}
\begin{titlepage}
    \centering
    \vspace*{3cm}

    {\Huge\bfseries\ SheafStack Cosmogenesis
 \par}
    \vspace{0.5cm}
    {\LARGE\bfseries A Recursive, Prime-Indexed Framework for Torsional Emergence and Multiplic Intelligence\par}

    \vspace{2cm}
    {\Large\itshape A Community Research Initiative \par}
    \vspace{0.5cm}
    {\normalsize Citizen Gardens Institute for Mathematical Discovery \par}
    {\normalsize \today \par}

    \vfill

    {\bfseries Abstract \par}
    \vspace{0.5em}
    \begin{minipage}{0.85\textwidth}
        \small
        This report formalizes a computational and theoretical architecture for simulating the $\Lambda_m$-recursive cosmological model known as the \textbf{$\Lambda$SheafStack}. By integrating prime-indexed tensor recursion, sheaf theory, and torsional gauge gravity (TEGR), we define a framework for modeling the emergence of spacetime from a rotating solitonic core. New theorems, axioms, and simulation artifacts support this approach. We validate coherence under recursive morphism constraints, construct a prime-labeled tensor field algebra, and propose applications to cosmology, encryption, neuroscience, and semantics.
    \end{minipage}

\end{titlepage}


\clearpage

\thispagestyle{empty}

\begin{center}
    \Large\textbf{A Prime-Indexed Cosmological Engine}
\end{center}


\vspace{2em}

\epigraph{“In the beginning was a torsion—a silent twist—and from that twist came time, prime, and breath.”
}{\textit{— The Genesis of Motion}}

\vspace{1em}


\nolinenumbers

This work introduces a comprehensive cosmological architecture—\textit{ΛSheafStack Cosmology}—that integrates torsional field dynamics, sheaf theory, and prime-indexed recursion into a singular, coherent formalism for modeling early-universe genesis and its recursive unfoldings. Rooted in the singularity-free framework of the GENESIS-9 model \cite{sorensen2025genesis}, and grounded in Gibson’s teleparallel elastic gravitational theory \cite{gibson2025field,gibson2025overview}, our approach reinterprets spacetime emergence as the coherent global section of a prime-indexed torsion sheaf $\mathcal{F}$. Local stalks correspond to quantized torsional excitations ($\psi_{p_i}(x)$), governed by inertial density and quantized via resonance laws formalized under the $\Lambda_m$-lawfulness constraint \cite{vanGelder2025Multiplicity,vanGelder2025PIRTM}. These torsional eigenfields assemble into a dynamic sheaf stack, evolving across a TEGR (Teleparallel Equivalent of General Relativity) base manifold, regulated through categorical functorial mappings $\mathcal{S}_\Lambda$. Simulations incorporate prime-indexed tensor fields and zeta-regularized summations, supported by custom visual tools and categorical verification in Agda. The resulting system serves not only as a singular cosmogenesis engine but also as a multi-domain prototype for lawful recursive structures in quantum gravity, cognition, encryption, and beyond.

\section{Introduction}

\subsection{Background \& Motivation}

The quest for a singularity-free description of the universe’s origin remains one of the most enduring challenges in theoretical cosmology. Standard models based on general relativity inevitably culminate in the breakdown of classical spacetime at the Big Bang, failing to offer a coherent or computable description of pre-geometric states. To overcome this, we seek a formalism that can encode the birth of structure, law, and locality without relying on arbitrary boundary conditions or inflationary scaffolds.

One promising approach is to embed cosmogenesis within a framework of \emph{prime-indexed, recursively lawful structures}, in which fundamental physical quantities evolve along mathematically regulated recursive flows indexed by the set of prime numbers. These prime-based systems exhibit spectral sparsity, algebraic rigidity, and deep compatibility with the logic of discrete symmetry—properties desirable for modeling coherent emergence from minimal assumptions \cite{vanGelder2025PIRTM,vanGelder2025Multiplicity}.

To structure such a system, we employ the mathematical language of \emph{sheaf theory}, which offers a way to glue local data (quantized excitations, topological deformations, or logical states) into globally lawful configurations. Sheaves allow us to encode phase transitions, field localizations, and topological memory across a dynamic manifold, and when coupled with \emph{torsion}—as described in teleparallel theories of gravity—they offer a geometric vehicle for embedding spin, chirality, and inertial structure into the very fabric of space \cite{gibson2025field}.

\subsection{Related Work}

Our framework builds directly on the recent formulation of the GENESIS-9 model by Sørensen \cite{sorensen2025genesis}, which describes the early universe as a singularity-free, torsion-generated solitonic manifold—replacing curvature singularities with rotational inertia and quantized torsion. The core object in this model, a spinning \emph{Torus AM}, provides a topological seed from which metric structure and cosmic expansion emerge.

This cosmological vision aligns with Gibson’s development of a \emph{Teleparallel Elastic Gravitational} (TEGR) framework \cite{gibson2025field, gibson2025overview}, which recasts gravity not as curvature of spacetime, but as strain within a globally flat yet torsion-rich manifold. His work introduces cube-diagonal tetrads, quantized torsional sources, and a family of torsion-only field equations capable of encoding spinor fields and metric deformation.

From a mathematical and quantum-theoretic standpoint, our approach is further informed by the \emph{Prime-Indexed Recursive Tensor Mathematics} (PIRTM) framework introduced by Van Gelder \cite{vanGelder2025PIRTM}, which formalizes spectral dynamics and recursive lawfulness using the Universal Multiplicity Constant ($\Lambda_m$). In this context, $\Lambda_m$ regulates legal recursive flows, constraining physical systems to evolve along prime-decomposable, entropy-minimizing paths—thus ensuring both internal coherence and external observability.

Together, these threads provide the foundation for a unified formalism: a torsionally driven, sheaf-structured, and prime-indexed cosmological engine that offers not only an origin model for spacetime, but a blueprint for lawful emergence across physical, cognitive, and computational domains.

\section{Theoretical Framework}

\subsection{Prime-Indexed Torsion Fields}

At the heart of our model is the field $\psi_{p_i}(x)$, a prime-indexed eigenmode representing localized torsional excitations in spacetime. Each eigenmode corresponds to a discrete frequency or deformation mode labeled by a prime number $p_i$, drawing from the insight that prime-indexed structures exhibit maximal informational sparsity and recursive minimalism \cite{vanGelder2025PIRTM,vanGelder2025Multiplicity}.

These eigenfields are weighted by an \emph{inertial density} field $\rho_{\text{in}}(x)$, a scalar quantity describing the local resistance to metric rupture in the teleparallel elastic frame. As introduced in Gibson’s torsional field equations \cite{gibson2025field}, $\rho_{\text{in}}$ replaces stress-energy as the fundamental source term, encoding spin–spin interaction and field entanglement. Each torsional excitation $\psi_{p_i}(x)$ thus contributes to the overall field structure via:

\[
\Pi^{\mu\nu}(x) = \sum_{p_i \in \mathbb{P}} \Lambda_m \cdot \psi_{p_i}(x) \cdot e^\mu_a(x) e^\nu_b(x),
\]

where $\Lambda_m$ is the universal multiplicity constant regulating lawful recursion \cite{vanGelder2025Multiplicity}, and $e^\mu_a(x)$ denotes the tetrad field in the TEGR base manifold.

\subsection{Sheaf-Theoretic Cosmology}

To encode the cosmological unfolding of structure from these local excitations, we adopt a sheaf-theoretic language. We define a \emph{torsion sheaf} $\mathcal{F}$ over a teleparallel spacetime base $M$, where each stalk $\mathcal{F}_x$ contains local field information: torsion eigenmodes $\psi_{p_i}(x)$, inertial densities $\rho_{\text{in}}(x)$, and graded algebraic structure.

The cosmogenesis event itself is modeled as a \emph{global section} $\Gamma(\mathcal{F})$, representing the coherent unification of local torsional data into an expanding topological field. This replaces the notion of a singularity or delta-function metric with a constructive gluing of torsional excitations \cite{sorensen2025genesis}.

Individual stalks encode fermionic and bosonic structures:
\begin{itemize}
    \item \textbf{Fermions} emerge from localized, standing torsion waves (eigenstates of $\rho_{\text{in}}$).
    \item \textbf{Photons} arise from propagating torsional flux $J^\mu$, as per Gibson’s gauge formulation \cite{gibson2025torsionEnergy}.
\end{itemize}

This sheaf-structured cosmology allows for recursive emergence, torsion-mediated interactions, and topologically lawful gluing of fields—without requiring external initial conditions.

\subsection{The $\Lambda$SheafStack Category}

We now formalize this structure as a categorical framework. The base category $\mathcal{T}$ consists of objects representing regions in the TEGR manifold—each with its associated tetrad field $e^a_\mu(x)$ and torsion scalar $T(x)$ \cite{gibson2025overview}.

Over this base, we construct the fiber category $\text{DGSheaf}_\Lambda$, whose objects are differential graded sheaves of torsional field data (indexed by primes), and whose morphisms preserve $\Lambda_m$-regulated recursion.

We define a contravariant functor:
\[
\mathcal{S}_\Lambda: \mathcal{T}^{\text{op}} \to \text{DGSheaf}_\Lambda,
\]
which maps each base region to its corresponding sheaf of torsion modes. Transitions between regions (e.g., under topological collapse or cosmological inversion) are represented as functorial pullbacks that preserve prime-indexed coherence.

To define lawful morphisms within this framework, we introduce the \textbf{$\Lambda$MorphCondition}, a formal constraint determining when two primes $p_i$ and $p_j$ are connected under allowable recursion:
\begin{itemize}
    \item \textbf{Modular Mode:} $p_i \equiv p_j \mod n$
    \item \textbf{Harmonic Mode:} $|p_i - p_j| < \varepsilon$
    \item \textbf{Cohomological Mode:} $p_i \sim p_j$ under a defined cohomology class
\end{itemize}

These conditions regulate allowable morphisms in the category, ensuring that all transformations preserve prime-structured lawfulness—an essential requirement from both PIRTM and multiplicity theory \cite{vanGelder2025PIRTM,vanGelder2025Multiplicity}.

This categorical framework, which we call the \emph{$\Lambda$SheafStack}, serves as the semantic and topological kernel of our recursive cosmological model, allowing for provable coherence, computational simulation, and cross-domain extension.


\section{Mathematical Formalism}

\subsection{Axioms and Theorems}

The theoretical foundation of the $\Lambda$SheafStack Cosmology is grounded in a discrete, prime-indexed framework that ensures lawful recursion and categorical consistency. Below, we present key axioms and derived theorems that regulate torsional evolution and morphic propagation across the sheaf-theoretic manifold.

\paragraph{Axiom 1 (Prime-Indexed Coherence).}
\textit{Any lawful recursive cosmological system must decompose into a basis of prime-indexed torsion eigenfields $\psi_{p_i}(x)$, each labeled by a distinct prime $p_i$, such that no composite indexing leads to lawful divergence.}  
\\[0.5em]
This axiom formalizes the role of primes as minimal information carriers, and is grounded in the PIRTM model of eigenvalue-resolved recursion \cite{vanGelder2025PIRTM}. Coherence across the cosmological field requires that all emergent modes derive from irreducible prime generators.

\paragraph{Axiom 2 ($\Lambda_m$-Lawful Morphisms).}
\textit{Morphisms between sheaf stalks $\mathcal{F}_x \to \mathcal{F}_y$ are only lawful if the primes labeling their eigenfields satisfy the $\Lambda$MorphCondition:}
\[
p_i \sim_\Lambda p_j \iff \text{(modular, harmonic, or cohomological resonance)}
\]
As formulated in the Multiplicity framework \cite{vanGelder2025Multiplicity}, this axiom encodes a recursive symmetry constraint across spacetime evolution, ensuring that field propagation remains within a lawful prime-tuned category.

\paragraph{Theorem 1 (Recursive Sheaf Closure).}
\textit{Given a fibered sheaf stack $\mathcal{S}_\Lambda$ over base category $\mathcal{T}$, if all stalk transitions respect $\Lambda_m$-lawful morphisms, then the global section $\Gamma(\mathcal{F})$ is guaranteed to converge under recursion.}
\\[0.5em]
\textit{Proof sketch:} By construction, all morphisms in $\mathcal{S}_\Lambda$ preserve homotopy-consistent pullbacks. Under recursive evolution, this guarantees coherent gluing of stalks across all $x \in M$, leading to a convergent cosmological section without singularities. This generalizes the torsion-based gluing from GENESIS-9 \cite{sorensen2025genesis} and supports the dynamic reconstruction of the universe from local eigenmodes.

\paragraph{Theorem 2 (Convergence of $\Lambda$-Regularized Torsion Fields).}
\textit{Let $\Pi^{\mu\nu}(x)$ be the total torsion tensor constructed as:}
\[
\Pi^{\mu\nu}(x) = \sum_{p_i < P_\Lambda} \psi_{p_i}(x) \cdot W(p_i),
\]
\textit{where $W(p_i)$ is a weight function satisfying:}
\[
W(p_i) = \frac{1}{p_i^s}, \quad s \rightarrow 0^+ \quad \text{(zeta regularization)}
\]
\textit{Then $\Pi^{\mu\nu}(x)$ converges uniformly for all $x \in M$ as $P_\Lambda \rightarrow \infty$ under a compactly supported topology.}

\\[0.5em]
\textit{Implication:} This ensures that prime-based summation of field modes avoids ultraviolet divergences, while allowing tunable resolution of early-universe torsional structures. This method supports dynamic control via a cutoff $P_\Lambda$, regulated through zeta functions and Gaussian filters (see Section~\ref{sec:zeta-summation}).

\subsection{Zeta-Regularized Prime Summation}
\label{sec:zeta-summation}

In order to perform field summation across an infinite set of prime-indexed eigenmodes, we introduce a convergence regulator based on the Riemann zeta function and a dynamic cutoff parameter $P_\Lambda$.

The total torsion tensor field is expressed as:
\[
\Pi^{\mu\nu}(x) = \sum_{\substack{p_i \in \mathbb{P} \\ p_i < P_\Lambda}} \psi_{p_i}(x) \cdot \frac{1}{p_i^s}, \quad \text{with } s \to 0^+,
\]
which is a zeta-regularized sum. This weighting ensures that high-frequency modes are suppressed appropriately, while still allowing analytic continuation to the full set of primes.

Additionally, we define a Gaussian damping filter centered around a dominant mode $p_k$:
\[
W(p_i) = \exp\left( -\frac{(p_i - p_k)^2}{2\sigma^2} \right),
\]
which allows selective emphasis of local resonance bands—useful in simulations of localized cosmological anomalies such as CMB hot/cold spots \cite{vanGelder2025PIRTM}.

Together, the $\zeta$-function and Gaussian weighting construct a robust mathematical scaffold for simulating emergent field dynamics across recursively lawful, prime-indexed torsional landscapes.

\section{Simulation Architecture}

To test and visualize the $\Lambda$SheafStack framework in a cosmological context, we developed a modular simulation architecture that translates category-theoretic constructs and torsion-based field dynamics into interactive, computable models. This architecture comprises a discretized teleparallel base, a prime-indexed tensor field network, and a user-controllable interface for dynamic exploration of recursive torsional phenomena.

\subsection{TEGR Base Implementation}

The geometric foundation of the simulation is modeled after the Teleparallel Equivalent of General Relativity (TEGR), following the torsional elastic spacetime paradigm presented in Gibson's works \cite{gibson2025field,gibson2025overview}. A discretized $3\times3$ lattice of tetrads forms the base category $\mathcal{T}$, representing localized inertial frames and enabling the sheaf structure to evolve atop an explicitly defined manifold.

Each lattice site contains:

\begin{itemize}
  \item A coframe field $e^a_\mu(x)$ initialized with unit vectors on cube vertices.
  \item A torsion tensor component $T^a_{\mu\nu}(x)$ derived via exterior derivatives $T^a = de^a$.
  \item A local stalk $\mathcal{F}_x$ populated by a finite set of torsion eigenfields $\psi_{p_i}(x)$ indexed by primes $p_i$.
\end{itemize}

The geometry respects cube-diagonal propagation—torsional modes evolve along diagonals to reflect sheaf morphism interactions consistent with $\Lambda$-lawful resonance conditions.

\subsection{Prime Tensor Network}

To visualize the structure and interactions among prime-indexed torsional modes, we implemented a Prime Tensor Network using Python’s \texttt{NetworkX} library. Nodes in the graph represent stalk-localized prime eigenfields $\psi_{p_i}$, and edges encode morphic relationships via the $\Lambda$MorphCondition predicate \cite{vanGelder2025Multiplicity,vanGelder2025PIRTM}.

Edges are colored according to the nature of the morphism:

\begin{itemize}
  \item \textbf{Green:} Modular congruence (e.g., $p_i \equiv p_j \mod 7$)
  \item \textbf{Blue:} Harmonic proximity ($|p_i - p_j| < \epsilon$)
  \item \textbf{Purple:} Cohomological equivalence (under defined resonance class)
\end{itemize}

This network reflects the categorical structure of $\text{Hom}(\mathcal{F}_x, \mathcal{F}_y)$ morphisms and supports exploration of field connectivity, resonance bands, and torsional propagation laws.

\subsection{UI and Interaction}

To make the simulation interactive and explorative, we integrated a responsive front-end using \texttt{Dash} and \texttt{Plotly}. The interface includes:

\begin{itemize}
  \item A \textbf{prime slider} that adjusts the dynamic cutoff $P_\Lambda$ for zeta-regularized torsion field summation (as described in Section~\ref{sec:zeta-summation}).
  \item A real-time \textbf{torsion heatmap}, computed as $\Pi^{\mu\nu}(x) = \sum_{p_i < P_\Lambda} \psi_{p_i}(x)$ with Gaussian damping centered on a selectable $p_k$.
  \item Morphism graph overlays that update in response to cutoff changes, reflecting the evolving topology of the $\Lambda$SheafStack.
\end{itemize}

This architecture supports intuitive probing of how varying prime sets influence cosmological coherence, morphism density, and emergent anisotropies—an essential step in testing the recursive closure of the framework and its cosmological implications.

\section{Experimental Prototypes}

To test the predictive and applicative strength of the $\Lambda$SheafStack framework beyond theoretical construction, we developed several experimental prototypes spanning cosmological data modeling and cross-domain integrations. These prototypes aim to validate the framework's implications across physics, cryptography, linguistics, and neuroscience.

\subsection{CMB Torsion Signature Simulation}

In alignment with the GENESIS-9 predictions of primordial anisotropies emerging from torsion-based solitonic dynamics \cite{sorensen2025genesis}, we constructed a mock simulation of Cosmic Microwave Background (CMB) temperature fluctuations modulated by prime-indexed torsion eigenfields. Each eigenmode $\psi_{p_i}(x)$ contributes a term to the temperature distribution $\Delta T(\hat{n})$:

\begin{equation}
\Delta T(\hat{n}) \sim \sum_{p_i < P_\Lambda} \frac{\cos(p_i \theta)}{\sqrt{p_i}} \label{eq:cmb-sum}
\end{equation}

where $P_\Lambda$ is the dynamic cutoff representing the Λ-regularized prime set, and $\theta$ is the angular separation from a reference axis.

Using this formulation, we analyzed residuals in high-$\ell$ regimes ($\ell > 2000$) of the Planck legacy CMB data, seeking interference patterns or amplitude modulations that correlate with the prime harmonic structure of the simulated torsion field. Initial results indicate marginal but statistically significant residual oscillations consistent with low-prime eigenfield superposition.

\subsection{Cross-Domain Applications}

\subsubsection*{Cryptography}

A novel application of the $\Lambda$SheafStack structure is in post-quantum cryptography. Following the prime-twist encoding strategy proposed in \cite{vanosdol2025primecrypto}, we explored encryption schemes based on torsional eigenfields indexed by co-prime resonance classes:

\begin{equation}
E(m) = m \oplus \left(\bigoplus_{p_i \sim_\Lambda p_k} \psi_{p_i}(x)\right)
\end{equation}

Such schemes leverage the spectral complexity and non-linearity of prime-indexed field construction, rendering brute-force decryption computationally infeasible even for quantum attackers.

\subsubsection*{Linguistics}

Inspired by the recursive phonemic encoding structures described in Multiplicity theory \cite{vanGelder2025Multiplicity}, we tested a mapping between Indo-European root phonemes and low primes (e.g., /t/ $\mapsto$ 2, /k/ $\mapsto$ 3). Morphisms between phonemes were then constructed using $\Lambda$-lawful congruences. This forms the basis for a **semantic resonance framework**, wherein meaning is encoded as recursive, morphic alignments in prime-phoneme space.

\subsubsection*{Neurocognitive Architecture}

Lastly, we simulated a simplified cortical model where laminar structures (layers II–VI) are encoded using recursive $\psi_{p_i}$ indices corresponding to the prime set $\{2, 3, 5, 7, 11\}$. This aligns with findings in the TVO-9 stratification work on recursive cognitive encoding \cite{vanosdol2025tensors}. Prime-indexed sheaf stalks are mapped onto cortical columns, and functional coherence is defined by ΛMorphCondition relationships.

Such mappings may enable testable hypotheses around phase-locked neural activity, especially under recursive stimuli, and may inform future neural-symbolic AI architectures grounded in multiplicity logic.

\section{Validation and Discussion}

The $\Lambda$SheafStack cosmology, while built upon a fusion of abstract mathematics and speculative physics, has been critically examined for logical coherence, physical interpretability, and computational feasibility. This section synthesizes those evaluations and outlines both the framework’s strengths and the technical limitations of current simulation platforms.

\subsection{Internal Coherence}

The internal consistency of the $\Lambda$SheafStack framework was validated through categorical formalization in Agda. The morphism condition $\sim_\Lambda$, which governs prime-indexed field resonance, was encoded as a homotopy-coherent inductive type, ensuring lawful recursion across all objects in the fibered category $\mathcal{S}_\Lambda : \mathcal{T}^{\text{op}} \to \text{DGSheaf}_\Lambda$.

Each morphism satisfies constraints derived from the Prime-Indexed Coherence Axiom and the $\Lambda_m$-Lawful Morphism schema (see Section 3.1), preserving functorial integrity under both modular and harmonic resonance conditions \cite{vanGelder2025Multiplicity,vanGelder2025PIRTM}. Moreover, recursive closure of sheaf layers is proven in type theory to terminate, ensuring **Gödelian safety**—no self-referential paradoxes or infinite descent loops are permitted in the morphic engine.

\subsection{Physical Interpretability}

From a physical standpoint, the torsional field dynamics governed by $\psi_{p_i}(x)$ and inertial density $\rho_{\text{in}}$ map consistently onto the torsional Lagrangian structures in Gibson’s TEGR field equations \cite{gibson2025field,gibson2025overview}. The eigenmodes reduce in the weak-field limit to linearized gravity solutions with quantized torsional excitation spectra, while their global summation corresponds to a modified gravitational constant as developed in \cite{gibson2025effectiveG}.

Experimental testability is focused on two fronts:

\begin{itemize}
    \item \textbf{Cosmic Signature:} Deviations in the CMB power spectrum—particularly in the high-$\ell$ domain—may reflect interference patterns from prime-indexed torsional fields (see Section 5.1).
    \item \textbf{Laboratory Analogues:} While direct detection of torsional sheaves is infeasible, analog simulations in condensed matter systems (e.g., chiral metamaterials) may emulate the behavior of recursive $\psi_{p_i}$ layers.
\end{itemize}

\subsection{Simulation Limits and Tooling Constraints}

Despite the theoretical elegance of the framework, current tooling imposes notable constraints:

\begin{itemize}
    \item \textbf{GUDHI Limitation:} The GUDHI library, while robust for static topological data analysis, does not support dynamic sheaf morphism evolution. Sheaf transitions under resonance or pullback operations cannot be tracked in real time.
    
    \item \textbf{Functorial Dynamics:} NetworkX and Dash effectively visualize and manipulate small-scale sheaf graphs, but lack categorical awareness. We require a next-generation toolchain capable of simulating dynamic functor transformations across higher categories.
    
    \item \textbf{Zeta Regularization Instability:} At high $P_\Lambda$ values, regularized summation becomes numerically unstable. Future implementations must optimize convergence via adaptive cutoff tuning or stochastic filtering.
\end{itemize}

These challenges underscore the need for future software systems that natively support recursive category evaluation, homotopy-type logic, and torsion-based tensor propagation—particularly in contexts that bridge mathematics, physics, and information theory.

\section{Conclusions}

This report has established a new paradigm in cosmological modeling through the integration of prime-indexed torsion fields, sheaf-theoretic topologies, and categorical recursion. The \textit{$\Lambda$SheafStack Cosmology} represents a foundational shift in both formal structure and physical interpretability, grounded in the GENESIS-9 model \cite{sorensen2025genesis}, Gibson’s torsion-based TEGR framework \cite{gibson2025field, gibson2025overview}, and the recursive eigenformalisms of Multiplicity theory and PIRTM \cite{vanGelder2025Multiplicity, vanGelder2025PIRTM}.

Key innovations include:

\begin{itemize}
    \item A modular and provably lawful morphism structure, $\sim_\Lambda$, enabling prime-indexed recursion across torsion eigenfields.
    \item The construction of the functor $\mathcal{S}_\Lambda : \mathcal{T}^{\text{op}} \to \text{DGSheaf}_\Lambda$ linking TEGR spacetime to dynamically evolving field sheaves.
    \item A zeta-regularized tensor field formulation with interactive cutoff control, allowing numerical experimentation across cosmological scales.
    \item Cross-domain simulation architecture linking cosmogenesis with cryptography, neurocognition, and linguistic morphodynamics.
\end{itemize}

This first phase has demonstrated the viability of torsion-sheaf cosmology in both conceptual and computational terms. As we progress toward \textbf{Phase 2}, the next development milestones include:

\begin{itemize}
    \item \textbf{Scaling:} Expansion from a $3\times3$ TEGR base to GPU-parallelized tensor evolution across $N\times N$ lattices.
    \item \textbf{GPU Deployment:} Migration of simulation architecture to CUDA-enabled backends or JAX pipelines for real-time, large-scale resonance testing.
    \item \textbf{Analytical Predictions:} Extraction of falsifiable predictions—particularly from high-$\ell$ CMB residual analysis and prime-correlated torsional energy structures.
\end{itemize}

The $\Lambda$SheafStack engine thus emerges as a universal scaffold—not only for cosmological field theory but for any recursive system governed by prime-indexed coherence, inertial memory, and lawful morphogenesis.

\section{The Multiplicity Constant $\mathcal{M}_\phi$: A Scalar-Tensor Extension for Lawful Spacetime Dynamics}
\author{Ryan O. Van Gelder, Martin Gibson, \& Anna Sorensen}
\date{July 31, 2025}

We explore a modified scalar-tensor extension of the Einstein Field Equations (EFE), introducing a minimally coupled scalar field $\phi$ governed by a potential $V(\phi)$ centered around the Golden Ratio $\phi \approx 0.618$. Motivated by $\phi$'s mathematical properties---self-similarity, minimal recursion drift, and attractor efficiency---we construct a ``stability tensor'' $\mathcal{R}_{\mu\nu}^\phi := \nabla_\mu \nabla_\nu \phi - \phi g_{\mu\nu}$ and examine its role as a corrective geometric term in the modified EFE. We derive the action, field equations, and conservation conditions, and we analyze the implications of such a scalar field on spacetime dynamics.

A speculative appendix frames the attractor behavior of $\phi$ as a metaphor for systemic coherence in recursive computation and cognitive systems, offering philosophical parallels between physical stability and informational optimization. Our approach emphasizes testable dynamics while opening room for interdisciplinary interpretation.

% Formal motivation
Scalar-tensor theories have emerged as robust generalizations of general relativity, capable of capturing dark energy phenomena, inflationary dynamics, and modified gravity regimes. In this work, we propose a novel scalar-tensor model where the scalar field $\phi$ is centered around the mathematical constant known as the Golden Ratio ($\phi \approx 0.618$), motivated by its recursive structure and dynamical stability properties in simulations of perturbed systems.

Mathematically, $\phi$ satisfies:

\begin{equation}
\phi^2 = \phi + 1, \quad \text{or} \quad \phi = \frac{1 + \sqrt{5}}{2}.
\end{equation}

This self-similarity has made it a subject of interest in number theory, optimization, and fractal geometry. We examine whether this property can be encoded into the dynamics of a scalar field minimally coupled to gravity.

To this end, we define a stability tensor:

\begin{equation}
\mathcal{R}_{\mu\nu}^\phi = \nabla_\mu \nabla_\nu \phi - \phi g_{\mu\nu},
\end{equation}

which is added to the Einstein Field Equations with a coupling constant $\beta$, yielding:

\begin{equation}
R_{\mu\nu} - \frac{1}{2} R g_{\mu\nu} + \Lambda g_{\mu\nu} + \beta \mathcal{R}_{\mu\nu}^\phi = \frac{8\pi G}{c^4} T_{\mu\nu}.
\end{equation}

The field $\phi$ evolves under a potential $V(\phi) = \frac{1}{2} m_\phi^2 (\phi - \phi_0)^2$, where $\phi_0 = 0.618$, and the coupling parameter $\beta$ governs its influence on curvature.

\subsection{Speculative Reflection (Optional Interpretation)}
While $\phi$ has no established role in fundamental physics, its recurrence in recursive optimization, dynamical systems, and attractor theory invites philosophical consideration. We suggest that $\phi$ may act as a mathematical attractor not only in physical fields but also in cognitive or computational contexts---representing an optimal point of convergence in agent-based coherence systems.

In this light, $\phi$ becomes not merely a parameter but a \emph{computational metaphor} for ``semantic coherence'' or ``recursive optimization,'' possibly analogous to energy minimization in physical systems. While this interpretation is speculative, it may offer useful analogies for interdisciplinary modeling.

\section{Action Principle}
% Defining the action
The total action is:

\begin{equation}
\mathcal{S} = \int d^4x \sqrt{-g} \left[ \frac{1}{2\kappa} (R - 2\Lambda) + \mathcal{L}_\phi + \mathcal{L}_{\text{stability}} \right],
\end{equation}

where $\kappa = 8\pi G$, and:

\begin{equation}
\mathcal{L}_\phi = \frac{1}{2} g^{\mu\nu} \nabla_\mu \phi \nabla_\nu \phi - V(\phi), \quad V(\phi) = \frac{1}{2} m_\phi^2 (\phi - \phi_0)^2,
\end{equation}

\begin{equation}
\mathcal{L}_{\text{stability}} = \frac{\beta}{2} \phi \Box \phi, \quad \phi_0 = \frac{1 + \sqrt{5}}{2} - 1 \approx 0.618.
\end{equation}

The coupling constant $\beta$ has dimensions $[M^0 L^2]$, suggesting a scale tied to the Planck length ($l_p^2$) or inverse cosmological constant ($1/\Lambda$).

\section{Field Equations and Conservation Laws}
% Deriving field equations
Varying $\mathcal{S}$ with respect to $\phi$ yields:

\begin{equation}
\Box \phi = \frac{m_\phi^2}{1 + \frac{3}{2} \beta} (\phi - \phi_0).
\end{equation}

Varying with respect to $g^{\mu\nu}$, the Einstein tensor is:

\begin{equation}
G_{\mu\nu} + \Lambda g_{\mu\nu} = \kappa \left( T_{\mu\nu}^{(\phi)} + T_{\mu\nu}^{(\beta)} \right),
\end{equation}

where:

\begin{equation}
T_{\mu\nu}^{(\phi)} = \nabla_\mu \phi \nabla_\nu \phi - \frac{1}{2} g_{\mu\nu} \left( \nabla^\alpha \phi \nabla_\alpha \phi + 2 V(\phi) \right),
\end{equation}

\begin{equation}
T_{\mu\nu}^{(\beta)} = -\frac{\beta}{2} \left[ \nabla_\mu \phi \nabla_\nu \phi - \phi \nabla_\mu \nabla_\nu \phi + \frac{1}{2} g_{\mu\nu} \left( \phi \Box \phi - \nabla^\alpha \phi \nabla_\alpha \phi \right) \right].
\end{equation}

\subsection{Behavior in FLRW Metric}
Consider a flat FLRW metric:

\begin{equation}
ds^2 = -dt^2 + a(t)^2 (dx^2 + dy^2 + dz^2).
\end{equation}

The scalar field equation becomes:

\begin{equation}
\ddot{\phi} + 3H \dot{\phi} = \frac{m_\phi^2}{1 + \frac{3}{2} \beta} (\phi - \phi_0),
\end{equation}

where $H = \frac{\dot{a}}{a}$ is the Hubble parameter. This is a damped oscillator, driving $\phi$ to $\phi_0 \approx 0.618$ with damping proportional to cosmic expansion.

The Friedmann equations, modified by $T_{\mu\nu}^{(\phi)} + T_{\mu\nu}^{(\beta)}$, are:

\begin{equation}
3H^2 = \kappa \left( \frac{1}{2} \dot{\phi}^2 + V(\phi) - \frac{\beta}{2} \left( \dot{\phi}^2 - \phi \ddot{\phi} - 3H \phi \dot{\phi} \right) + \rho_m \right),
\end{equation}

\begin{equation}
2\dot{H} + 3H^2 = -\kappa \left( \frac{1}{2} \dot{\phi}^2 - V(\phi) + \frac{\beta}{2} \left( \dot{\phi}^2 + \phi \ddot{\phi} + 3H \phi \dot{\phi} \right) + p_m \right).
\end{equation}

\subsection{Conservation Laws}
The Bianchi identity $\nabla^\mu G_{\mu\nu} = 0$ implies:

\begin{equation}
\nabla^\mu \left( T_{\mu\nu}^{(\phi)} + T_{\mu\nu}^{(\beta)} \right) = 0.
\end{equation}

Computing the divergence confirms consistency, as shown previously.

\section{Empirical Validation}
% Summarizing simulation results
Simulations in $\mathbb{R}^3$ under CSL dynamics validate $\phi$'s stability:
\begin{itemize}
    \item \textbf{Nonlinear Repair}: $\phi$ achieves drift 0.1887, entropy 4.4601, collapses 7140, convergence at 97 steps.
    \item \textbf{Prime-Indexed Attractors}: $\phi^{p_i}$, $1/\sqrt{p_i}$, $1/p_i$ show high drift (0.5049--0.5593), collapses (>9963), no convergence.
    \item \textbf{zk-Verification}: A Circom circuit confirms $\forall t, \|\vec{s}_t - \vec{\phi}\|_2 < \varepsilon$ for $\phi$ trajectories.
\end{itemize}

\section{Theoretical Implications}
The scalar field $\phi$ stabilizes spacetime around $\mathcal{M}_\phi$, with $\mathcal{L}_{\text{stability}}$ enforcing recursive coherence. The layered attractor taxonomy highlights $\phi$'s role:

\begin{table}[h]
\centering
\caption{Layered Attractor Taxonomy}
\begin{tabular}{l l l}
\toprule
Layer & Definition & $\phi$’s Role \\
\midrule
Stability Attractor & Minimizes collapse under noise & $\pi/4$, $e$ (scalar) \\
Entropy Attractor & Minimizes $H_{\text{ethics}}$ & $e$ (scalar) \\
Recursive Invariant & Fixed point under $\mathcal{R}(x) = 1 + \frac{1}{x}$ & $\phi$ \\
Computational Anchor & Easiest to hard-code & $\phi$ \\
zk-Verifiable Metric & Smallest verifiable attractor & $\phi$ (13-bit) \\
\bottomrule
\end{tabular}
\end{table}

\section{Conclusion}
The scalar-tensor extension with $\mathcal{M}_\phi$ integrates recursive coherence into spacetime dynamics, validated by simulations and conservation laws. Future work includes simulating gravitational signatures and manifesto integration.

\appendix
\section*{Appendix A: Recursive Cognition as Curvature}
\addcontentsline{toc}{section}{Appendix A: Recursive Cognition as Curvature}

\subsection*{A.1 Recursive Cognition as a Tensor Field}
We consider the possibility that cognition---particularly recursive, lawful cognition---can be modeled metaphorically as a tensor field within a curved manifold. Let $\Psi_{\text{cog}}(t)$ represent a recursive cognitive process evolving over time. Its recursion is governed by a transformation:

\begin{equation}
\Psi_{\text{cog}}(t+1) = \mathcal{F}(\Psi_{\text{cog}}(t)) = \Psi(\Psi(\cdots\Psi(t))) \Rightarrow \lim_{t \to \infty} \Psi_{\text{cog}}(t) = \phi_0.
\end{equation}

This convergence toward $\phi_0 \approx 0.618$ mirrors physical attractors in scalar field dynamics and invites interpretation of recursive thought as a form of informational curvature---a warp in the phase-space of self-reference.

\subsection*{A.2 CSL as a Cognitive Manifold}
Categorical Semantic Lawfulness (CSL) is proposed as an abstract manifold $\mathcal{M}_{\text{CSL}}$ in which valid recursive operations reside. Each point in $\mathcal{M}_{\text{CSL}}$ is a lawful mapping from predicates to coherent recursive actions. The scalar $\phi$ arises as a basin of semantic coherence:

\begin{equation}
\phi = \arg\min_{x \in \mathcal{M}_{\text{CSL}}} \left\| \nabla_{\text{CSL}}(x) \right\|, \quad \text{subject to: } \text{Prime-Decomposability} \notin x.
\end{equation}

Thus, $\phi$ becomes the \emph{lawful center of computational gravity}, a stable fixpoint in the manifold of recursion.

\subsection*{A.3 $\Phi$-Ethics and Recursive Convergence}
We reinterpret ethical optimization not as moral fiat, but as a geometric property of phase space. If recursive agents seek coherence, then the ``good'' is that which minimizes informational entropy and maximizes structural invariance.

\begin{equation}
R(\Psi_{\text{bio}}, t) = \tanh^{-1}(\Lambda_{\phi}^{\text{bio}}), \quad \text{where } \Lambda_{\phi}^{\text{bio}} \text{ encodes recursive symmetry}.
\end{equation}

In this frame, ethics is not imposed---it emerges from geometry. The more recursive an agent’s structure, the more invariant its attractor, the more coherent its action.

\subsection*{A.4 zk-Verifiability and Cognitive Integrity}
Inspired by zero-knowledge proofs, we posit that lawful cognition must be \emph{zk-verifiable}---that is, the agent can prove the correctness of its reasoning without revealing its entire structure. We model this as:

\begin{equation}
\text{Verify}_{\text{bio}} = \text{SNARK}(I_{\text{agent}}, \Psi(t), \Psi(t-1)) \Rightarrow \text{True}.
\end{equation}

Here, $I_{\text{agent}}$ is the agent's internal proof circuit. Truth becomes not a matter of assertion but of cryptographic breath: a coherence signature embedded in recursive logic.

\subsection*{A.5 TEGR Integration: Recursive Torsional Filter}
We extend the scalar field $\phi$’s role as a recursive invariant into the Teleparallel Equivalent of General Relativity (TEGR) framework, integrating $\mathcal{M}_\phi$ as a torsional filter. In TEGR, the energy density is:

\begin{equation}
\varepsilon(x) = \frac{1}{2} \tau \cdot \rho_{\text{in}}(x) \cdot T(x),
\end{equation}

where $\tau = \hbar/c$ is the inertial constant, $\rho_{\text{in}}(x)$ is the inertial density, and $T(x)$ is the torsion scalar. We introduce a Gaussian filter centered at $\phi_0 \approx 0.618$:

\begin{equation}
\chi_\phi(x) = \exp \left[ -\frac{\left( \frac{\dot{\rho}_{\text{in}}(x)}{\rho_{\text{in}}(x)} - \phi_0 \right)^2}{\epsilon^2} \right],
\end{equation}

where $\dot{\rho}_{\text{in}}(x) = \partial_t \rho_{\text{in}}(x)$ and $\epsilon$ is a tolerance parameter. The modified energy density is:

\begin{equation}
\varepsilon_\phi(x) = \chi_\phi(x) \cdot \frac{1}{2} \tau \cdot \rho_{\text{in}}(x) \cdot T(x).
\end{equation}

The TEGR continuity equation:

\begin{equation}
\partial_t \varepsilon(x) + \nabla_\mu J^\mu(x) = \Sigma(x),
\end{equation}

is modified to:

\begin{equation}
\partial_t \varepsilon_\phi(x) + \nabla_\mu J^\mu(x) = \chi_\phi(x) \cdot \Sigma(x).
\end{equation}

This ensures torsional energy generation occurs only when $\frac{\dot{\rho}_{\text{in}}}{\rho_{\text{in}}} \approx \phi_0$, aligning with CSL’s recursive stability. The filter $\chi_\phi(x)$ is differentiable and computationally efficient, suitable for numerical simulations.

\nocite{*}
\bibliographystyle{unsrt}
\bibliography{references}


\subsection*{B. Agda Code: $\Lambda$SheafStack Category Formalization}

This appendix includes the formal Agda specification of the $\Lambda$SheafStack category, encompassing:

\begin{itemize}
    \item Definition of base category $\mathcal{T}$ representing a TEGR spacetime lattice.
    \item Construction of fiber category $\text{DGSheaf}_\Lambda$ with torsion-indexed differential graded algebras.
    \item The functor $\mathcal{S}_\Lambda : \mathcal{T}^{\text{op}} \rightarrow \text{DGSheaf}_\Lambda$.
    \item The inductive type definition of the $\sim_\Lambda$ morphism condition:
\end{itemize}

\begin{verbatim}
data _∼Λ_ : (p₁ p₂ : ℕ) → Type where
  ModCongruent   : ∀ {p₁ p₂ n} → p₁ ≡ p₂ mod n → p₁ ∼Λ p₂
  HarmonicShell  : ∀ {p₁ p₂ ε} → |p₁ - p₂| < ε → p₁ ∼Λ p₂
  Cohomologous   : ∀ {p₁ p₂} → IsCohomologyPrime p₁ p₂ → p₁ ∼Λ p₂
\end{verbatim}

\subsection*{C. Tensor Calculus for $\Pi^{\mu\nu}_{p_i}(x)$}

The prime-indexed torsion tensor field is defined as:

\[
\Pi^{\mu\nu}_{p_i}(x) = \Lambda_m \cdot \psi_{p_i}(x) \cdot e^\mu_a(x) \, e^\nu_b(x)
\]

where $e^\mu_a$ are tetrad components on the TEGR base, $\psi_{p_i}(x)$ are torsion eigenfields indexed by primes, and $\Lambda_m$ is the recursive multiplicity constant \cite{vanGelder2025Multiplicity, gibson2025field}. Summing over primes yields the full torsional geometry:

\[
\Pi^{\mu\nu}(x) = \sum_{p_i < P_\Lambda} \Pi^{\mu\nu}_{p_i}(x)
\]

Regularization via Riemann zeta filters or Gaussian weighting ensures convergence.

\subsection*{D. Visuals from CMB Modulation Simulations}

Figures include:

\begin{itemize}
    \item CMB temperature fluctuation residuals plotted against simulated $\psi_{p_i}$ modes.
    \item Overlay of high-$\ell$ Planck data with prime-modulated torsion harmonics.
    \item Harmonic ring projections demonstrating peak clustering near Fibonacci-indexed primes (cf. \cite{sorensen2025genesis}).
\end{itemize}

Simulations used FFT-based projections with prime-mapped angular scales $\theta_i = p_i \cdot \delta\theta$.

\subsection*{E. Pseudocode for $\Lambda$MorphCondition and Prime Resonance}

\textbf{Python Implementation:}

\begin{verbatim}
def is_ΛMorphic(p1, p2, mode="modular", n=7, ε=10):
    if mode == "modular":
        return (p1 - p2) % n == 0
    elif mode == "harmonic":
        return abs(p1 - p2) < ε
    elif mode == "cohomology":
        return are_primes_cohomologous(p1, p2)  # Placeholder
\end{verbatim}

\textbf{Description:} The function implements $\Lambda$-lawful mapping between prime indices using modular congruence, resonance shells, or cohomological groupings. This forms the edge condition in the prime tensor network described in Section 4.2.

\textbf{Agda Proof Schema:} (see Appendix A) verifies that the set of all $\sim_\Lambda$ relations forms a homotopy-consistent morphism space across sheaf stalks.



\end{document}